\chapter{State-space dynamics, neural manifolds, and rotations}\label{ch:dynamics}
\begin{readergoal}Connect dynamical evolution, latent inference, and geometric structure while keeping descriptive and mechanistic claims separate.\end{readergoal}
\section{A state is more than a snapshot}
A discrete linear dynamical system has $z_{t+1}=Az_t+Bu_t+w_t$ and observations $y_t=Cz_t+v_t$. The hidden state $z_t$ evolves according to a transition rule while observations mix it with noise. If the model is linear Gaussian, Kalman filtering estimates the posterior state sequentially. Smoothing uses future as well as past observations and is not equivalent to an online estimate.

\section{Continuous dynamics and stability}
A continuous linear system obeys $\dot z=Az+Bu(t)$. The eigenvalues of $A$ govern local behavior: negative real parts give asymptotic decay for unforced linear systems, while complex eigenvalues generate damped rotational components. More generally $\dot z=f(z,u,t)$ can represent nonlinear dynamics. Stable equilibria and attractors are properties of the chosen model; inferring them from noisy, partial observations requires care.

\section{jPCA and rotational population trajectories}
Principal-component analysis selects directions of large variance, not necessarily directions displaying dynamics. The jPCA approach first reduces dimensionality, then fits an approximately skew-symmetric linear flow for the derivatives of the reduced trajectories, highlighting rotations \cite{Churchland2012}. The existence of a rotational projection is a description of population activity under chosen preprocessing; attributing that rotation to a particular recurrent mechanism requires additional evidence.

\section{Neural manifolds and coordinate nonuniqueness}
A low-dimensional manifold is a set locally parameterized by a smaller number of coordinates than its ambient dimension. If observations obey $y=g(z)$ with $z\in\mathbb R^d$, many coordinate systems for $z$ can describe the same observed trajectories. Thus latent axes may not have unique biological labels. A fitted low-dimensional model can be useful for decoding even when its coordinates are not uniquely identifiable. Reviews of movement-related manifolds emphasize dynamical and task constraints \cite{Gallego2017}.

\section{Nonlinear latent models}
Sequential autoencoders and recurrent neural networks can infer nonlinear latent dynamics from single trials, under training objectives and architecture assumptions. Models such as LFADS use latent generative dynamics for population spike data. Flexible models can also overfit. Appropriate controls include held-out likelihood, behavior decoding, simulation-based identifiability tests, and stability across recording sessions.

\begin{workedexample}[A two-dimensional rotation]
Let $A=\begin{bmatrix}0&-\omega\\\omega&0\end{bmatrix}$. Then $\dot z=Az$ has solutions $z_1(t)=R\cos(\omega t+\phi)$ and $z_2(t)=R\sin(\omega t+\phi)$. The trajectory is a circle of radius $R$. A change of orthonormal coordinates rotates the picture but leaves its essential frequency unchanged.
\end{workedexample}
\begin{exercises}
\item Determine eigenvalues of the rotation matrix above.
\item Describe a dynamical system with a stable fixed point.
\item Explain why PCA dimensions and dynamical dimensions need not coincide.
\item Name two empirical checks required before claiming a neural manifold controls behavior.
\end{exercises}